% GENERATED FILE. Edit canonical lessons, then run npm run generate:books.
% canonical-source-hash: fnv1a64:ab5a44ca989d213c
% canonical-lessons: TA-C475-cikiccai, TA-C475-paricoti, TA-C475-nati, TA-C475-sken, TA-C475-pattiyam

\chapter{At the Doctor}
\label{ch:ta-at-the-doctor}
\hlchaptermodality{475}

\begin{chapteropening}
\textbf{By the end of this chapter:} \emph{I can talk about treatment, check-ups, the pulse, scans and diets.}

Complete the last lesson of chapter 475: I can talk about treatment, check-ups, the pulse, scans and diets.
\end{chapteropening}

\section[\texorpdfstring{\ta{சிகிச்சை} (cikiccai)}{சிகிச்சை (cikiccai)}]{\ta{சிகிச்சை} (cikiccai) — treatment; therapy}

\label{lesson:TA-C475-cikiccai}



\begin{quote}
Before the new one: say the Tamil for to hiccup, then the Tamil for to chafe.
\end{quote}

\subsection*{You'll want to know: \ta{சிகிச்சை}}
\textbf{\ta{சிகிச்சை}} — \emph{cikiccai} — \textquotedblleft{}treatment; therapy\textquotedblright{}.

\begin{quote}
\ta{அவருக்கு} \ta{மருத்துவமனையில்} \ta{சிகிச்சை} \ta{அளித்தார்கள்}.
\end{quote}

\begin{quote}
\emph{avarukku maruttuvamaṉaiyil cikiccai aḷittārkaḷ.}
\end{quote}

\begin{quote}
They gave him treatment in the hospital.
\end{quote}

\ta{சிகிச்சை} \ta{பெறு}, to receive treatment.

\begin{cousinweb}
\textbf{\ta{சிகிச்சை}} (\emph{cikiccai}) is borrowed from Sanskrit \emph{cikitsā}, \textquotedblleft{}medical attendance, therapeutics\textquotedblright{} (Monier-Williams).
\end{cousinweb}

\begin{grammarlens}[title={What you've built}]
One more word for describing what you lived and felt.
\end{grammarlens}

\subsection*{Your turn}
\begin{itemize}
\raggedright
  \item \emph{Say it:} \emph{cikiccai}
  \item \emph{Say it:} \emph{cikiccai}, once more
  \item \emph{Say it:} \emph{avarukku maruttuvamaṉaiyil cikiccai aḷittārkaḷ.} — They gave him treatment in the hospital.
  \item \emph{Recall:} say the Tamil for to hiccup, then the Tamil for to chafe, then say \emph{cikiccai} again
\end{itemize}

\begin{tcolorbox}[breakable,colback=teal!4,colframe=teal!35!black,title={Before you move on}]
What does \ta{சிகிச்சை} mean? (\textquotedblleft{}Treatment; therapy\textquotedblright{}.) Say it once more.
\end{tcolorbox}
\section[\texorpdfstring{\ta{பரிசோதி} (paricōti)}{பரிசோதி (paricōti)}]{\ta{பரிசோதி} (paricōti) — to examine; to test}

\label{lesson:TA-C475-paricoti}



\begin{quote}
Before the new one: say the Tamil for to chafe, then the Tamil for treatment.
\end{quote}

\subsection*{You'll want to know: \ta{பரிசோதி}}
\textbf{\ta{பரிசோதி}} — \emph{paricōti} — \textquotedblleft{}to examine; to test\textquotedblright{}.

\begin{quote}
\ta{மருத்துவர்} \ta{என்} \ta{கண்ணைப்} \ta{பரிசோதித்தார்}.
\end{quote}

\begin{quote}
\emph{maruttuvar eṉ kaṇṇaip paricōtittār.}
\end{quote}

\begin{quote}
The doctor examined my eye.
\end{quote}

Past (I …): \ta{பரிசோதித்தேன்}, \emph{paricōtittēṉ}.

\ta{பரிசோதனை}, a medical test.

\begin{cousinweb}
\textbf{\ta{பரிசோதி}} (\emph{paricōti}) is borrowed from Sanskrit: the prefix \emph{pari-} with the stem of \emph{śōdhayati}, \textquotedblleft{}purifies, clears\textquotedblright{} (CDIAL 12630).
\end{cousinweb}

\begin{grammarlens}[title={What you've built}]
One more word for describing what you lived and felt.
\end{grammarlens}

\subsection*{Your turn}
\begin{itemize}
\raggedright
  \item \emph{Say it:} \emph{paricōti}
  \item \emph{Say it:} \emph{paricōti}, once more
  \item \emph{Say it:} \emph{maruttuvar eṉ kaṇṇaip paricōtittār.} — The doctor examined my eye.
  \item \emph{Recall:} say the Tamil for to chafe, then the Tamil for treatment, then say \emph{paricōti} again
\end{itemize}

\begin{tcolorbox}[breakable,colback=teal!4,colframe=teal!35!black,title={Before you move on}]
What does \ta{பரிசோதி} mean? (\textquotedblleft{}To examine; to test\textquotedblright{}.) Say it once more.
\end{tcolorbox}
\section[\texorpdfstring{\ta{நாடி} (nāṭi)}{நாடி (nāṭi)}]{\ta{நாடி} (nāṭi) — the pulse}

\label{lesson:TA-C475-nati}



\begin{quote}
Before the new one: say the Tamil for treatment, then the Tamil for to examine.
\end{quote}

\subsection*{You'll want to know: \ta{நாடி}}
\textbf{\ta{நாடி}} — \emph{nāṭi} — \textquotedblleft{}the pulse\textquotedblright{}.

\begin{quote}
\ta{மருத்துவர்} \ta{என்} \ta{நாடியைப்} \ta{பார்த்தார்}.
\end{quote}

\begin{quote}
\emph{maruttuvar eṉ nāṭiyaip pārttār.}
\end{quote}

\begin{quote}
The doctor felt my pulse.
\end{quote}

\ta{நாடி} \ta{பார்}, to take the pulse.

\begin{cousinweb}
\textbf{\ta{நாடி}} (\emph{nāṭi}) is borrowed from Sanskrit \emph{nāḍī}, \textquotedblleft{}a tube, a tubular organ\textquotedblright{} (CDIAL 7047): the vessel where the pulse is felt.
\end{cousinweb}

\begin{grammarlens}[title={What you've built}]
One more word for describing what you lived and felt.
\end{grammarlens}

\subsection*{Your turn}
\begin{itemize}
\raggedright
  \item \emph{Say it:} \emph{nāṭi}
  \item \emph{Say it:} \emph{nāṭi}, once more
  \item \emph{Say it:} \emph{maruttuvar eṉ nāṭiyaip pārttār.} — The doctor felt my pulse.
  \item \emph{Recall:} say the Tamil for treatment, then the Tamil for to examine, then say \emph{nāṭi} again
\end{itemize}

\begin{tcolorbox}[breakable,colback=teal!4,colframe=teal!35!black,title={Before you move on}]
What does \ta{நாடி} mean? (\textquotedblleft{}The pulse\textquotedblright{}.) Say it once more.
\end{tcolorbox}
\section[\texorpdfstring{\ta{ஸ்கேன்} (skēṉ)}{ஸ்கேன் (skēṉ)}]{\ta{ஸ்கேன்} (skēṉ) — a scan}

\label{lesson:TA-C475-sken}



\begin{quote}
Before the new one: say the Tamil for to examine, then the Tamil for the pulse.
\end{quote}

\subsection*{You'll want to know: \ta{ஸ்கேன்}}
\textbf{\ta{ஸ்கேன்}} — \emph{skēṉ} — \textquotedblleft{}a scan\textquotedblright{}.

\begin{quote}
\ta{மருத்துவர்} \ta{வயிற்றுக்கு} \ta{ஒரு} \ta{ஸ்கேன்} \ta{எடுக்கச்} \ta{சொன்னார்}.
\end{quote}

\begin{quote}
\emph{maruttuvar vayiṟṟukku oru skēṉ eṭukkac coṉṉār.}
\end{quote}

\begin{quote}
The doctor asked for a scan of the stomach.
\end{quote}

\ta{ஸ்கேன்} \ta{எடு}, to have a scan taken.

\begin{cousinweb}
\textbf{\ta{ஸ்கேன்}} (\emph{skēṉ}) is a loan from English \emph{scan}, written with the grantha letter \ta{ஸ}.
\end{cousinweb}

\begin{grammarlens}[title={What you've built}]
One more word for describing what you lived and felt.
\end{grammarlens}

\subsection*{Your turn}
\begin{itemize}
\raggedright
  \item \emph{Say it:} \emph{skēṉ}
  \item \emph{Say it:} \emph{skēṉ}, once more
  \item \emph{Say it:} \emph{maruttuvar vayiṟṟukku oru skēṉ eṭukkac coṉṉār.} — The doctor asked for a scan of the stomach.
  \item \emph{Recall:} say the Tamil for to examine, then the Tamil for the pulse, then say \emph{skēṉ} again
\end{itemize}

\begin{tcolorbox}[breakable,colback=teal!4,colframe=teal!35!black,title={Before you move on}]
What does \ta{ஸ்கேன்} mean? (\textquotedblleft{}A scan\textquotedblright{}.) Say it once more.
\end{tcolorbox}
\section[\texorpdfstring{\ta{பத்தியம்} (pattiyam)}{பத்தியம் (pattiyam)}]{\ta{பத்தியம்} (pattiyam) — a restricted diet}

\label{lesson:TA-C475-pattiyam}



\begin{quote}
Before the new one: say the Tamil for the pulse, then the Tamil for a scan.
\end{quote}

\subsection*{You'll want to know: \ta{பத்தியம்}}
\textbf{\ta{பத்தியம்}} — \emph{pattiyam} — \textquotedblleft{}a restricted diet\textquotedblright{}.

\begin{quote}
\ta{காய்ச்சலுக்குப்} \ta{பிறகு} \ta{பாட்டி} \ta{ஒரு} \ta{வாரம்} \ta{பத்தியம்} \ta{இருந்தாள்}.
\end{quote}

\begin{quote}
\emph{kāyccalukkup piṟaku pāṭṭi oru vāram pattiyam iruntāḷ.}
\end{quote}

\begin{quote}
After the fever grandmother kept to a strict diet for a week.
\end{quote}

\ta{பத்தியம்} \ta{இரு}, to keep to a diet, often in traditional medicine.

\begin{cousinweb}
\textbf{\ta{பத்தியம்}} (\emph{pattiyam}) is borrowed from Sanskrit \emph{pathya}, \textquotedblleft{}suitable, wholesome\textquotedblright{} (CDIAL 7746).
\end{cousinweb}

\begin{grammarlens}[title={What you've built}]
One more word for describing what you lived and felt.
\end{grammarlens}

\subsection*{Your turn}
\begin{itemize}
\raggedright
  \item \emph{Say it:} \emph{pattiyam}
  \item \emph{Say it:} \emph{pattiyam}, once more
  \item \emph{Say it:} \emph{kāyccalukkup piṟaku pāṭṭi oru vāram pattiyam iruntāḷ.} — After the fever grandmother kept to a strict diet for a week.
  \item \emph{Recall:} say the Tamil for the pulse, then the Tamil for a scan, then say \emph{pattiyam} again
\end{itemize}

\begin{tcolorbox}[breakable,colback=teal!4,colframe=teal!35!black,title={Before you move on}]
What does \ta{பத்தியம்} mean? (\textquotedblleft{}A restricted diet\textquotedblright{}.) Say it once more.
\end{tcolorbox}
